\section{Preliminares}
La unidad mínima de información en un sistema de computación cuántica es el cúbit, denotado por un vector unitario en un espacio de Hilbert de dimensión 2 que se denota $\Q$.
Múltiples cúbits pueden tomarse conjuntamente de forma que se considera un elemento del espacio de Hilbert de $n$ cúbits, denotado por $\Q_n$.
Este caso puede verse como el producto tensorial de los $n$ espacios de Hilbert, $\Q_n=\Q^{\otimes n}=\Q\otimes\dots\otimes\Q$ correspondientes a cada cúbit del sistema, de modo que la dimensión del espacio resultante es $2^n$.

Sean $\N$, $\Z$, $\R$ y $\C$ los conjuntos de los números naturales, enteros, reales y complejos, respectivamente.
Como el $0$ no es un número natural, sea $\N^*$ el conjunto de los números naturales incluyendo el cero.
La notación $\N_n$ y $\N^*_n$ se usa para denotar el subconjunto de $\N$ y $\N^*$ respectivamente que contiene los números naturales menores o iguales que $n$.
En teoría de conjuntos, para cualquier conjunto finito $S$, su número de elementos $|S|$ se llama la cardinalidad de $S$.

Sea $\B$ el grupo $\set{0,1}$ con la suma binaria $\oplus$, y la delta de Kronecker es la función $\delta:\N^*\times\N^*\longrightarrow\B$ definida por $\delta(k,l)=\delta_{k,l}=1$ si $k=l$ y $0$ en otro caso.
Para un conjunto $\Gamma\subseteq\N^*_n$, el valor de $\delta_{k,\Gamma}$ es $1$ si $k\in\Gamma$ y $0$ en otro caso.
